% =====================================
% Le modèle et sa dérivation
% =====================================

\subsection{Modèle de Malthus}
\begin{frame}
    \begin{itemize}
        \item On considère une population homogène, isolée, qui vit dans un environnement stable avec des ressources illimitées. On note $P(t)$ le nombre total d'individu à l'instant $t \in \bbR_+$. \pause
        %\item Dans ces conditions, les variations de la population vont uniquement dépendre des taux de fertilité et de mortalité, notés respectivement $\beta$ et $\mu$, et supposés constants. \pause
        \item On observe la variation de population sur un intervalle de la forme $[t,t+h]$. On a 
        \begin{equation*}
            P(t+h) = P(t) + \text{naissances} - \text{décès}
        \end{equation*} \pause
        \item Sur l'intervalle $[t,t+h]$, le nombre de naissance est $\beta P(t)h$ et le nombre de décès est $\mu P(t) h$. On substitue dans l'équation précédente et on obtient
        \begin{equation*}
            P(t+h)=P(t)+\beta P(t)h - \mu P(t)h
        \end{equation*} \pause
        \item En divisant par $h$ et en faisant tendre $h$ vers $0$, on obtient finalement \textcolor{astral}{\textit{l'équation de Malthus}}
        \begin{equation*}
            P'(t) = (\beta - \mu)P(t)
        \end{equation*} %\pause
        %\item La solution de cet EDO est donnée par $P(t) = P(0)\ex^{(\beta - \mu)t}$.
    \end{itemize}
\end{frame}

\subsection{Modèle de McKendrick-Lotka}
\begin{frame}
    \begin{itemize}
        \item On va décrire l'évolution de la population par une fonction donnant la \textcolor{astral}{\textit{densité d'individu d'âge $a$ à l'instant $t$}}
        \begin{equation*}
            x(a,t) \geq 0, \ \ \ \ \ a \in [0,a_{\max}], \ t \geq 0
        \end{equation*}
        %où $a_{\max}$ est l'âge maximal de la population, supposé fini. 
        \pause
        \item On introduit la fonction $\beta$ qui représente le \textcolor{astral}{\textit{taux de fertilité de la population}} en fonction de l'âge. On peut alors considérer le nombre total de nouveau né à l'instant $t$
        \begin{equation*}
            B(t):= \int_{0}^{a_{\max}}\beta(a)x(a,t)d a
        \end{equation*} \pause
        \item On introduit à présent la fonction $\mu$ qui représente le \textcolor{astral}{\textit{taux de mortalité}} de la population en fonction de l'âge. On définit ensuite la \textcolor{astral}{\textit{probabilité de survie}}
        \begin{equation*}
            \Pi (a) := \exp\bigg(-\int_{0}^a \mu(\sigma)d \sigma\bigg)
        \end{equation*}
        et on pose $\Pi(a_{\max}) = 0$. \pause
        \item On définit de plus la \textcolor{astral}{\textit{fonction de maternité}} pour tout $a \in [0,a_{\max}]$, $K(a) := \beta(a)\Pi(a)$. \pause 
        \item Enfin, on définit le \textcolor{astral}{\textit{nombre de reproduction}}
        \begin{equation*}
            R := \int_0^{a_{\max}} \beta(a)\Pi(a)d a = \int_0^{a_{\max}} K(a)d a
        \end{equation*}
    \end{itemize}
\end{frame}

\begin{frame}
    \begin{itemize}
        \item On peut à présent définir le modèle de McKendrick-Lotka. Comme pour le modèle de Malthus, on étudie la variation de la population sur $[t,t+h]$
        \begin{equation*}
            x(a+h,t+h) = x(a,t) - \text{décès}
        \end{equation*} \pause 
        \item Le taux de décès entre $t$ et $t+h$ est donné par
        \begin{equation*}
            \int_0^{h}\mu(a+s)x(a+s,t+s)ds = \int_t^{t+h}\mu(a+s-t)x(a+s-t,s)d s
        \end{equation*} \pause  
        %\item On pose $G(t) := \displaystyle{\int_0^t}\mu(a+s-t)x(a+s-t,s)ds$, d'où
        %\begin{equation*}
        %    x(a+h,t+h)-x(a,t+h)+\big(x(a,t+h) - x(a,t)\big) = - \big(G(t+h)-G(t)\big)
        %\end{equation*} \pause
        \item On divise par $h$ et on fait tendre $h$ vers $0$, on obtient alors \textcolor{astral}{\textit{l'équation de McKendrick-Lotka}}
        \begin{equation*}
            \partial_tx(a,t)+\partial_ax(a,t) = -\mu(a)x(a,t)
        \end{equation*} \pause
        \item On obtient le \textcolor{astral}{\textit{modèle de McKendrick-Lotka}}
        \begin{equation*}
            \left\{\begin{array}{l}
            \partial_t x(a,t) + \partial_a x(a,t) = -\mu(a)x(a,t) \\
            x(0,t) = B(t) = \displaystyle{\int_0^{a_{\max}}} \beta (a)x(a,t)d a \\
            x(a,0) = x_0(a)
            \end{array}\right.
        \end{equation*}
    \end{itemize}
\end{frame}